\chapter{Operational Risk and Rogue Trading}\label{ch:rc:operational-risk-and-rogue-trading}

On 26 February 1995 the London bank Barings collapsed. Its trader in Singapore had run up losses in an account numbered 88888 that London never saw: more
than 20 million pounds by the end of 1993, more than 200 million by the end of 1994, and 827 million by
the collapse. He was in charge of both the trading and the back office that recorded it; the bank's
internal auditors had recommended in 1994 that he should not be; London had sent him more than 300
million pounds to lend to clients who had posted about 31 million of collateral, and nobody asked
why. None of the models of the previous chapters would have measured this risk: the positions were
hidden from them. This chapter is about the risk that sits outside the models: \omterm{def:rc:operational-risk-and-rogue-trading:oprisk}{operational risk}, and
its most expensive form in trading, unauthorised positions concealed by the people who run them.

\section{Operational risk and its measurement}

\begin{definition}[Operational risk, operational loss event]\label{def:rc:operational-risk-and-rogue-trading:oprisk}
\emph{Operational risk}\index{operational risk} is the risk of loss resulting from inadequate or
failed internal processes, people and systems, or from external events; it includes legal risk and
excludes strategic and reputational risk. An \emph{operational loss event}\index{operational loss event}
is an occurrence of such a loss, recorded with its date, amount, business line and cause.
\end{definition}

Operational losses are rare and heavy-tailed: most are small errors, a few are catastrophic. Banks
collect them in loss databases, map them to causes and controls, and hold capital against them.

\begin{method}[The Basel standardised approach]\label{met:rc:operational-risk-and-rogue-trading:sa}
Compute the business indicator (BI), a measure of the bank's size from its income statement. The
business indicator component (BIC) applies marginal coefficients of 12\% up to EUR~1 billion of BI, 15\% from
1 to 30 billion and 18\% above. Scale it by the internal loss multiplier
$\mathrm{ILM} = \ln\bigl(e-1+(\mathrm{LC}/\mathrm{BIC})^{0.8}\bigr)$, where the loss component LC is 15 times the average annual
operational loss of the last ten years. Capital is $\mathrm{BIC}\times\mathrm{ILM}$.
\end{method}

\begin{example}[Capital for operational risk]\label{ex:rc:operational-risk-and-rogue-trading:capital}
A bank with a BI of EUR~35 billion has a BIC of $0.12+29\times0.15+5\times0.18 = 5.37$ billion. With average
annual losses of 0.5 billion, its loss component is 7.5 billion, its ILM 1.107 and its capital 5.94 billion;
with 0.25 billion of losses, the ILM is 0.904 and the capital 4.85 billion.
\end{example}

\section{The documented cases}

\begin{definition}[Rogue trading, fictitious trade]\label{def:rc:operational-risk-and-rogue-trading:rogue}
\emph{Rogue trading}\index{rogue trading} is trading beyond a trader's authority, in size or in kind,
concealed from the firm. A \emph{fictitious trade}\index{fictitious trade} is a booking with no economic
reality, entered to offset or disguise a real position in the firm's records.
\end{definition}

\begin{dated}{2026-09}{Six cases}\label{dat:rc:operational-risk-and-rogue-trading:cases}
Barings (1995): losses of GBP~827 million concealed in account 88888 by a trader who also ran the
back office (the Chancellor's statement on the Board of Banking Supervision's report). Daiwa (1995):
a New York trader lost about USD~1.1 billion on Treasury securities over eleven years, covering the
losses by selling securities held in custody and falsifying records; he also ran the custody department
(Federal Reserve order). Sumitomo (1995--96): the US CFTC found that after large losses from
speculative trading, the company's principal copper trader engaged in a scheme to manipulate copper
prices; the company paid USD~150 million, and in September 1996 put its losses at USD~2.6 billion.
Allfirst (1997--2002): losses of USD~691.2 million concealed by fictitious foreign exchange options that
appeared to hedge real positions (the parent bank's annual report). Soci\'et\'e G\'en\'erale (2008): a long
position of EUR~49 billion on index futures, built from 2 to 18 January, cost EUR~6.4 billion to unwind,
EUR~4.9 billion net of 2007 gains; it was hidden by 947 \omterm{def:rc:operational-risk-and-rogue-trading:rogue}{fictitious trades} cancelled before confirmation
(the bank's inspection report). UBS (2011): losses of USD~2.3 billion on an ETF desk, concealed by late
bookings, unmatched internal trades and false trades with deferred settlement; the UK regulator fined
the bank GBP~29.7 million.
\end{dated}

\section{The control failures they share}

\begin{definition}[Segregation of duties, block leave]\label{def:rc:operational-risk-and-rogue-trading:sod}
\emph{Segregation of duties}\index{segregation of duties} separates the people who trade from those who
book, confirm, settle, value and report the trades. \emph{Block leave}\index{block leave} requires each
trader to be away from the desk, with no access to its systems, for a consecutive period each year, so that
positions that need daily attention to stay hidden are exposed.
\end{definition}

The documented failures share a pattern (\cref{fig:rc:operational-risk-and-rogue-trading:chain}): a
trader who controls the recording or custody of his own trades (Barings, Daiwa); internal counterparties
and trades that no one matches or confirms (UBS), or confirmations the trader forges (Allfirst); \omterm{def:rc:operational-risk-and-rogue-trading:rogue}{fictitious
trades} with distant value dates, cancelled before they are ever confirmed (Soci\'et\'e G\'en\'erale); funding
that no one reconciles with the positions it pays for (Barings). Each is a link in the chain of controls
that the trade was supposed to pass through, and each is broken by one person.

\begin{omfigure}
\begin{tikzpicture}[font=\small,box/.style={draw,rounded corners,minimum width=2.3cm,minimum height=0.95cm,align=center,text width=2.2cm},
  note/.style={omThm,align=center,font=\scriptsize,text width=2.6cm,anchor=north}]
\node[box,fill=omThm!10] (t) at (0,0) {trade};
\node[box] (b) at (3,0) {booking};
\node[box] (c) at (6,0) {confirmation};
\node[box] (s) at (9,0) {settlement and cash};
\node[box] (p) at (12,0) {P\&L explain and IPV};
\draw[->,thick] (t) -- (b);
\draw[->,thick] (b) -- (c);
\draw[->,thick] (c) -- (s);
\draw[->,thick] (s) -- (p);
\node[note] at (3,-0.75) {late bookings; no off-market filter (UBS)};
\node[note] at (6,-0.75) {unmatched internal trades (UBS)};
\node[note] at (9,-0.75) {deferred settlement (UBS); funding not reconciled (Barings)};
\node[note] at (12,-0.75) {losses in a hidden account (Barings)};
\node[omDef,align=center,font=\footnotesize] at (6,1.0) {one person controlling trading and back office breaks every link (Barings)};
\end{tikzpicture}

\omcaption{The chain of controls a trade passes through, and where the documented cases broke it.
Each link is independent evidence that the position exists and is what it claims to be. Schematic;
cases from the Chancellor's statement on Barings (1995) and the FSA's 2012 notice to UBS.}\label{fig:rc:operational-risk-and-rogue-trading:chain}
\end{omfigure}

\section{Detecting fictitious and hidden trades}

Surveillance turns the pattern into rules on the trade records: a trade booked long after it was
executed; cancellations and amendments around reporting dates, when positions are measured; prices far
from the market; internal trades with no matching trade on the other book; settlement dates far in the
future, for trades never meant to settle.

\omcode{code/firm/tradecontrol/firm_tradecontrol.py}{36}{55}{Five surveillance rules on a blotter of trades.}\label{lst:rc:operational-risk-and-rogue-trading:rules}

\begin{example}[Rules on a synthetic blotter]\label{ex:rc:operational-risk-and-rogue-trading:rules}
A synthetic blotter holds 2\,000 genuine trades, with the noise of real operations (4\% booked late, a few
off-market fills, internal trades of which 2\% lack a mirror, cancellations at random times), and 60 trades
of a concealment scheme. Alone, each rule catches 23\% to 52\% of the scheme's trades, with false-alarm rates
from 0.25\% (unmatched internal trades) to 4.4\% (late bookings). Flagging any trade that breaks one rule
catches 88.3\% of the scheme at a false-alarm rate of 7.1\%, 142 genuine trades to review; requiring two
rules catches 61.7\% with a single false alarm (\cref{fig:rc:operational-risk-and-rogue-trading:rules}).
\end{example}

\begin{omfigure}
\begin{tikzpicture}
\begin{axis}[width=11.5cm,height=5.6cm,ybar,bar width=8pt,
  symbolic x coords={late booking,cancel near report,off-market,internal unmatched,deferred settle,any rule,two rules},xtick=data,
  x tick label style={font=\footnotesize,rotate=35,anchor=north east},
  ylabel={\% of trades flagged},
  tick label style={font=\small},label style={font=\small},
  ymin=0,ymax=100,grid=major,grid style={gray!25},
  legend columns=2,legend style={font=\footnotesize,at={(0.5,1.03)},anchor=south,draw=none,fill=none,/tikz/every even column/.append style={column sep=8pt}},
  table/col sep=comma]
\addplot[fill=omDef!60,draw=omDef] table[x=rule,y=hit]{figdata/rates-credit-risk/28-operational-risk-and-rogue-trading/rules.csv};
\addplot[fill=omThm!60,draw=omThm] table[x=rule,y=false]{figdata/rates-credit-risk/28-operational-risk-and-rogue-trading/rules.csv};
\legend{scheme trades caught (hit rate),genuine trades flagged (false alarms)}
\end{axis}
\end{tikzpicture}

\omcaption{Hit and false-alarm rates of the surveillance rules on the synthetic blotter, alone and
combined. No single rule catches most of the scheme; combining them catches most of it, at the cost of
alerts that must be investigated. Data: synthetic; the chapter's tutorial.}\label{fig:rc:operational-risk-and-rogue-trading:rules}
\end{omfigure}

Alerts are only as good as their investigation: an alert closed on the explanation of the person it
concerns controls nothing. An append-only audit log, in which each entry carries the hash of the previous one,
makes the record of trades and of their investigation tamper-evident: an edited entry breaks the chain.

\section{Tutorial: surveillance on a blotter}\label{tut:rc:operational-risk-and-rogue-trading:tutorial}

\begin{tutorial}
\textbf{Goal.} Run the five rules on a synthetic blotter with a planted scheme, measure hit and false-alarm
rates, keep a tamper-evident log, and compute operational-risk capital. \textbf{End state:} the numbers of
\cref{ex:rc:operational-risk-and-rogue-trading:capital,ex:rc:operational-risk-and-rogue-trading:rules}
and the chart.
\begin{tutsteps}
\item \textbf{Blotter}: \texttt{blotter()}, genuine and scheme trades with labels.
\item \textbf{Rules}: \texttt{run\_rules}, \texttt{evaluate} with one and two rules required.
\item \textbf{Audit}: \texttt{audit\_demo()}: the chain verifies, then fails after an edit.
\item \textbf{Capital}: \texttt{op\_capital(35, 0.5)}; \texttt{fig\_rc\_tradecontrol.py}.
\end{tutsteps}
\textbf{What to change next.} Weight the rules by their false-alarm rates; add a rule on profits too
large for the desk's risk (a Sharpe ratio above a limit).
\end{tutorial}

\section{Build: trade control}\label{bld:rc:operational-risk-and-rogue-trading:tradecontrol}

\begin{build}
\bfield{Purpose} The firm's surveillance of its own trades: rules on every booking, alert scoring, and a
tamper-evident audit log for trades, amendments and investigations.
\bfield{Interface} \texttt{Trade}; \texttt{late\_booking}, \texttt{cancel\_amend\_near}, \texttt{off\_market},
\texttt{internal\_unmatched}, \texttt{deferred\_settlement}; \texttt{run\_rules}, \texttt{evaluate};
\texttt{AuditLog} with \texttt{append} and \texttt{verify}; \texttt{trade\_event}.
\bfield{Rules} Rules are independent and cheap; scores add; alerts are closed only with evidence recorded
in the log.
\bfield{Acceptance tests} \texttt{code/firm/tradecontrol/tests/}: each rule flags its planted case and not
the clean one; evaluation arithmetic; the audit chain verifies and detects an edit.
\bfield{Stretch} Rules on P\&L and cash (profits against risk, margin against positions); trader-level
behaviour baselines; links to confirmations and settlement data; block-leave monitoring.
\end{build}

\begin{omsources}
\item House of Commons, \emph{Hansard}, 18 July 1995, statement on the Board of Banking Supervision's
report on Barings.
\item Financial Services Authority, Final Notice to UBS AG, 25 November 2012.
\item Commodity Futures Trading Commission, release 4144-98, 11 May 1998.
\item Board of Governors of the Federal Reserve System, order against Daiwa Bank, 2 October 1995.
\item Allied Irish Banks, annual report on Form 20-F for 2001 (Allfirst).
\item Soci\'et\'e G\'en\'erale, General Inspection Department, \emph{Mission Green}, report of 20 May 2008.
\item Basel Committee on Banking Supervision, \emph{Basel~III: Finalising post-crisis reforms}, December
2017 (operational risk).
\end{omsources}

\section{Exercises}

\begin{exercise}[$\star$]\label{exo:rc:operational-risk-and-rogue-trading:1}
Compute the BIC of a bank with a business indicator of EUR~10 billion.
\end{exercise}

\begin{exercise}[$\star$]\label{exo:rc:operational-risk-and-rogue-trading:2}
Why is the ILM equal to one when the loss component equals the BIC?
\end{exercise}

\begin{exercise}[$\star$]\label{exo:rc:operational-risk-and-rogue-trading:3}
Which rule would have flagged the UBS scheme's deferred-settlement ETF trades?
\end{exercise}

\begin{exercise}[$\star\star$]\label{exo:rc:operational-risk-and-rogue-trading:4}
Why is lending GBP~300 million against GBP~31 million of collateral a control signal?
\end{exercise}

\begin{exercise}[$\star\star$]\label{exo:rc:operational-risk-and-rogue-trading:5}
Why does requiring two rules cut false alarms far more than hits?
\end{exercise}

\begin{exercise}[$\star\star$]\label{exo:rc:operational-risk-and-rogue-trading:6}
How does \omterm{def:rc:operational-risk-and-rogue-trading:sod}{block leave} expose a concealed position?
\end{exercise}

\begin{exercise}[$\star\star\star$]\label{exo:rc:operational-risk-and-rogue-trading:7}
\emph{Coding.} Verify that the audit log detects an edit to one trade's price, and explain why rewriting
the whole chain after an edit would still be detected if the latest hash is stored elsewhere.
\end{exercise}

\begin{exercise}[$\star\star\star$]\label{exo:rc:operational-risk-and-rogue-trading:8}
\emph{Find the flaw.} ``Our trader's P\&L is steady and high with low VaR; he is our best, and we should raise
his limits.''
\end{exercise}

\section{Problem: January 2008}

\begin{problem}[{Weekend problem --- sizing a hidden position}]\label{pb:rc:operational-risk-and-rogue-trading:1}
Société Générale's inspection report found that a trader built a long position of EUR~49 billion on
index futures between 2 and 18 January 2008, discovered on the 20th; unwinding it from 21 to 23 January
cost EUR~6.4 billion, or EUR~4.9 billion net of the 1.5 billion his positions had gained in 2007.

\medskip\noindent\textbf{Part I --- The arithmetic.}
\begin{enumerate}
\item Give the average adverse move implied by the loss and the position.
\item What position would the same loss imply at an average adverse move of 5\%?
\item Why did the unwind itself move the market against the bank?
\item Assuming an annual volatility of 25\% for the shares, what one-day 99\% VaR would the position
have shown?
\item Why did the firm's VaR not show it?
\end{enumerate}

\medskip\noindent\textbf{Part II --- The controls.}
\begin{enumerate}[resume]
\item What kind of \omterm{def:rc:operational-risk-and-rogue-trading:rogue}{fictitious trades} hide a large directional position?
\item Which of the chapter's rules would catch them?
\item What would confirmations with counterparties have revealed?
\item Which cash or margin signal should have appeared?
\item What role does \omterm{def:rc:operational-risk-and-rogue-trading:sod}{block leave} play?
\end{enumerate}

\medskip\noindent\textbf{Part III --- Detection in practice.}
\begin{enumerate}[resume]
\item Give the hit and false-alarm rates of the combined rules on the chapter's blotter.
\item How many alerts a year would a desk of that size generate, and who investigates them?
\item Why do alerts get closed without being resolved?
\item How does the audit log help an investigation after the fact?
\item What would you add to the rules for a futures desk?
\end{enumerate}

\medskip\noindent\textbf{Part IV --- Judgement.}
\begin{enumerate}[resume]
\item Why do rogue traders usually start by hiding losses, not by seeking gains?
\item Which control is the first line that would have caught it?
\item How should operational-risk capital reflect such events?
\item State the \emph{named result}: the average move implied by the reported loss and position, and the
first control in the chain that would have caught it.
\item In one sentence: what makes a control work?
\end{enumerate}
\end{problem}

\section{Interview questions}

\begin{interviewq}[$\star$ \iqroles{risk, bank}]\label{iq:rc:operational-risk-and-rogue-trading:1}
Define \omterm{def:rc:operational-risk-and-rogue-trading:oprisk}{operational risk}. How is it different from market and credit risk?
\end{interviewq}

\begin{interviewq}[$\star\star$ \iqroles{risk}]\label{iq:rc:operational-risk-and-rogue-trading:2}
What controls prevent \omterm{def:rc:operational-risk-and-rogue-trading:rogue}{rogue trading}? Which ones failed at Barings?
\end{interviewq}

\begin{interviewq}[$\star\star$ \iqroles{developer}]\label{iq:rc:operational-risk-and-rogue-trading:3}
Design a surveillance system for \omterm{def:rc:operational-risk-and-rogue-trading:rogue}{fictitious trades}.
\end{interviewq}

\begin{interviewq}[$\star\star$ \iqroles{trader}]\label{iq:rc:operational-risk-and-rogue-trading:4}
A colleague's P\&L is steady, large and unexplained by his positions. What do you do?
\end{interviewq}

\begin{interviewq}[$\star\star\star$ \iqroles{bank, risk}]\label{iq:rc:operational-risk-and-rogue-trading:5}
How is operational-risk capital computed under the Basel standardised approach, and what are its weaknesses?
\end{interviewq}

\begin{interviewq}[$\star\star\star$ \iqroles{risk}]\label{iq:rc:operational-risk-and-rogue-trading:6}
Why do surveillance alerts fail to stop unauthorised trading, and how would you fix it?
\end{interviewq}
